\chapter{Experiments and Results}
\label{chap:results}

% Sources: outputs/*/README.md, experiments/*.py,
%          docs/Design Principles.md

% Every result in this chapter concerns the Integrating paper~\cite{Benchasattabuse2025Integrating}: unless the Bridging paper~\cite{Benchasattabuse2026Bridging} is named explicitly, every mention of ``the source paper'' below refers to the Integrating paper, the source of both the Fig.\,4 hand-picked schedule (Section~\ref{sec:exp:baselines}) and the Fig.\,5 fidelity benchmarks (Section~\ref{sec:exp:validation}) this chapter's comparisons are built around. This chapter uses \emph{the paper schedule} for that Fig.\,4 example throughout, to avoid ambiguity with the schedule \Gls{dag} formalism of Chapter~\ref{chap:model}.

\TOFIX{repeating the paramters and sweeps part? fix oth paragraphs}{All experiments use the beam search ($w_{\text{beam}}=25$, Section~\ref{sec:alg:beam}) together with the fixed schedule families of Section~\ref{sec:alg:bruteforce}, unless stated otherwise. The reference configuration is $\mathcal{N}_{\mathrm{ref}}$ with $N=10$, $e_d=0.01$, $F_{\min}=0.9$, and $E_{\max}=100$, matching the source paper's main experimental setting. Unless the Bridging paper~\cite{Benchasattabuse2026Bridging} is named explicitly, ``the source paper'' in this chapter refers to the Integrating paper~\cite{Benchasattabuse2025Integrating}, whose Fig.\,4 schedule and Fig.\,5 fidelity results provide the main comparison points.}

This chapter first establishes that the evaluator reproduces the source paper's physical results, then compares the paper schedule with schedules found by search. It next examines how those results change under non-ideal and heterogeneous network conditions, followed by cost-quality tradeoffs and scaling with chain length. The chapter closes with timing sensitivity; additional runtime, beam-width, and sweep results are collected in Appendix~\ref{app:sweeps}.

\section{Experimental Setup}
\label{sec:exp:setup}

% Source: experiments/*.py parameter tables

\paragraph{Parameters and sweeps}
\label{sec:exp:params}

Unless stated otherwise, every experiment starts from the reference configuration $\mathcal{N}_{\mathrm{ref}}$ (Section~\ref{sec:model:network}) and changes only the parameters specified for that experiment; all other parameters remain fixed. Table~\ref{tab:params} lists the parameters varied in this chapter.

\begin{table}[ht]
    \centering
    \small
    \begin{tabular}{p{5cm}lllp{3cm}}
        \toprule
        Parameter & Symbol & Reference ($\mathcal{N}_{\mathrm{ref}}$) & Swept ranges & Where varied \\
        \midrule
        Hop count & $N$ & 10 & 10--28 & Section~\ref{sec:exp:budget} \\
        Hop length & $\ell_i$ & 2~km & \{1, 2, 3, 4\}~km & Section~\ref{sec:exp:heterogeneous} \\
        Tree-encoding branching vector & $\vec b^{(i)}$ & $(16,14,1)$ & -- & not varied \\
        BSM arm count & $k_i$ & 18 & 6, 18 & Section~\ref{sec:exp:netgen} \\
        Inner-qubit error rates & $\{p_X^{\mathrm{in}},p_Z^{\mathrm{in}}\}_i$ & 0 & [0, $3\!\times\!10^{-3}$] & Section~\ref{sec:exp:netgen} \\
        Photon survival probability & $\eta_i$ & 0.912 & $10^{-0.2 \ell_i/10}$ & Section~\ref{sec:exp:heterogeneous} \\
        Outer-photon depolarizing rate & $e_d$ & 0.01 & $[0,\,0.01]$ & Section~\ref{sec:exp:edsweep} \\
        Memory dephasing rate & $\gamma$ & 0 & $[0,\,10^5]$ & Section~\ref{sec:exp:timing} \\
        Propagation speed & $c$ & $2\!\times\!10^5$~km/s & -- & not varied \\
        Emission delay & $\tau_{\mathrm{emit}}$ & 0 & $[10^{-4},\,1]$ & Section~\ref{sec:exp:timing} \\
        Generation-cost budget & $E_{\max}$ & 100 & 50--280 & Sections~\ref{sec:exp:headline}--\ref{sec:exp:budget} \\
        \bottomrule
    \end{tabular}

    \caption{Experimental parameters. The reference column gives the nominal value used unless a parameter is varied in the indicated experiment. All parameters not listed as varying follow the Integrating paper's~\cite{Benchasattabuse2025Integrating} reference configuration $\mathcal{N}_{\mathrm{ref}}$ (Section~\ref{sec:model:network}). The tree branching factor is not varied because it does not affect the objectives $F$, $C$, or $R$ introduced in Chapter~\ref{chap:model}. Uniform hop-length scaling likewise leaves relative schedule comparisons unchanged; heterogeneous hop lengths are considered separately in Section~\ref{sec:exp:heterogeneous}.}
    \label{tab:params}
\end{table}

\paragraph{Search method runtime}
\label{sec:exp:runtime}

Both the exact dynamic program (Section~\ref{sec:alg:dp}) and the beam search, using beam width 25 and with the same-span pumping move enabled (Section~\ref{sec:alg:pumping}), complete within the tested range $N\le18$ without reaching the 300~s timeout imposed on each configuration. Table~\ref{tab:runtimes} reports the corresponding wall-clock times under the paper's $E_{\max}=10N$ budget.

\begin{table}[ht]
\centering
\caption{Wall-clock runtime and returned score for the exact dynamic program and beam search at representative chain lengths. Unless otherwise stated, each run uses the paper budget $E_{\max}=10N$, fidelity threshold $F_{\min}=0.9$, beam width $w_{\text{beam}}=25$, same-span pumping enabled (Section~\ref{sec:alg:pumping}), a 300~s per-point timeout, and a 3~GiB memory limit. The score is the value of the quantity maximized by the active objective (Section~\ref{sec:alg:beam}); under the default rate-maximizing objective used here, the score is the returned schedule rate $R(\Sigma)$. \textsuperscript{$\dagger$}At $N=18$, the exact dynamic program completes within the timeout but finds no feasible candidate, so its score is $-\infty$, whereas the beam search returns a feasible candidate. \textsuperscript{$\ddagger$}At $N=6$, the exact dynamic program returns a lower score than the beam search. These cases occur because same-span pumping is itself beam-limited in the implementation, so the pumping-enabled exact dynamic program is not a guaranteed upper bound on the beam search (Section~\ref{sec:alg:pumping}).}
\label{tab:runtimes}
\small
\begin{tabular}{rrllrr}
\toprule
$N$ & $E_{\max}$ & Exact DP (s) & Beam search (s) & Exact DP score & Beam search score \\
\midrule
2  & 20  & 0.10 & 0.10 & 50\,000 & 50\,000 \\
6  & 60  & 3.82 & 3.34 & 14\,372\textsuperscript{$\ddagger$} & 15\,391 \\
10 & 100 & 20.43 & 16.50 & 6\,196 & 6\,196 \\
14 & 140 & 85.37 & 66.93 & 3\,394 & 3\,394 \\
18 & 180 & 203.68\textsuperscript{$\dagger$} & 153.52 & $-\infty$ & 1\,214 \\
\bottomrule
\end{tabular}
\end{table}

\paragraph{Baselines}
\label{sec:exp:baselines}

Three fixed non-adaptive baselines, introduced in Section~\ref{sec:alg:bruteforce}, are used throughout: the \emph{raw chain}, with no purification; the \emph{uniform link-level baseline}, which applies the same copy count and circuit sequence at every hop; and the \emph{paper schedule}, the source paper's hand-picked Fig.\,4 example~\cite{Benchasattabuse2025Integrating}. As stated earlier, the latter is a demonstration of scheduling flexibility rather than an optimized design. The uniform link-level baseline and the paper schedule therefore represent distinct fixed strategies and are labelled separately below. The difference becomes particularly relevant under non-ideal inner-qubit noise, where the paper schedule can become infeasible while the uniform baseline remains feasible.

\paragraph{Metrics}
\label{sec:exp:metrics}

Fidelity $F$ is the Bell-pair fidelity after all purification, backbone, and correction operations, i.e.\ the $w$ component of the error vector at the schedule root (Section~\ref{sec:model:objectives}). Rate is $R=P_{\mathrm{success}}/T_{\mathrm{attempt}}$, where $P_{\mathrm{success}}$ is the probability that a complete schedule attempt succeeds and $T_{\mathrm{attempt}}$ is its deterministic wall-clock duration. Resource cost $C$ counts the number of \Gls{hrgs} copies, or equivalently \texttt{GenNode} leaves, consumed per trial. Two comparison modes recur throughout this chapter: \emph{matched-cost}, in which the optimizer is constrained to $C=100$, and \emph{budget-relaxed}, in which it may choose any schedule satisfying $C\le E_{\max}$.

\section{Physics Validation}
\label{sec:exp:validation}

% Source: experiments/fig5_fidelity_vs_noise.py,
%          outputs/headline_experiment_n10/README.md

The evaluator of Section~\ref{sec:alg:evaluator} reproduces the source paper's Fig.\,5~\cite{Benchasattabuse2025Integrating} to four significant figures. At $e_d=0.01$, the raw chain, heralded end-node pumping, and optimistic end-node pumping give $F=0.82$, $0.92$, and $0.93$, respectively, matching the paper's reported values. This agreement validates the physical model used by the search tiers, but does not establish search optimality. Appendix~\ref{app:validation} gives the corresponding comparison across the swept range of $e_d$.

\section{Optimizer versus the Paper Schedule}
\label{sec:exp:comparison}

This section first compares the search results with the source paper's schedule at the reference configuration, then examines the effect of non-ideal inner-qubit noise and outer-photon noise.

\subsection{Headline Comparison}
\label{sec:exp:headline}

% Source: outputs/headline_experiment_n10/README.md,
%          outputs/sweep_network_sensitivity/README.md

At $N=10$ and $e_d=0.01$, the optimizer improves on the paper schedule at both matched cost and under a relaxed budget. Table~\ref{tab:headline} reports the four variants. At matched cost, the optimizer returns $F=0.9168$ and $R=4158.14$ at $C=100$, compared with $F=0.9295$ and $R=4055.92$ for the paper schedule, giving a $2.5\%$ rate improvement while remaining above the fidelity floor. With the budget relaxed and same-span pumping enabled (Section~\ref{sec:alg:pumping}), the optimizer returns $C=60$, $F=0.9063$, and $R=6195.95$, using $60\%$ of the paper schedule's cost while increasing rate by $52.8\%$. With pumping disabled, the search instead returns $C=50$, $F=0.9047$, and $R=6713.18$. Both search settings are retained in the comparison because the pumping extension changes the beam-limited search space and is studied separately in Section~\ref{sec:alg:pumping}.

\begin{table}[ht]
\centering
\caption{Headline comparison at the reference configuration $\mathcal{N}_{\mathrm{ref}}$ ($N=10$, $e_d=0.01$, $F_{\min}=0.9$, $E_{\max}=100$). The paper schedule is the hand-picked Fig.\,4 construction from the Integrating paper~\cite{Benchasattabuse2025Integrating}. Matched-cost search fixes $C=100$; budget-relaxed search allows any schedule with $C\le100$. Rate is $R=P_{\mathrm{success}}/T_{\mathrm{attempt}}$.}
\label{tab:headline}
\small
\begin{tabular}{lcccl}
\hline
Schedule & $C$ & $F$ & Rate & Note \\
\hline
Paper schedule & 100 & 0.93 & 4055.92 & hand-picked \\
Matched-cost & 100 & 0.92 & 4158.14 & +2.5\% rate \\
Budget-relaxed & 60 & 0.91 & 6195.95 & +52.8\% rate, 60\% of cost \\
Budget-relaxed, pumping disabled & 50 & 0.90 & 6713.18 & lower cost still \\
\hline
\end{tabular}
\end{table}

\subsection{Network Generalization}
\label{sec:exp:netgen}

% Source: outputs/sweep_network_sensitivity/README.md

The headline comparison uses the source paper's idealized network configuration. Introducing non-zero inner-qubit error changes which schedules remain feasible. At the same generation cost, the paper schedule can fall below the $F\ge0.9$ threshold, while a search rerun under the changed physical parameters finds a different schedule that remains feasible. Table~\ref{tab:netgen} and Figure~\ref{fig:netgen} compare the paper schedule with the optimizer under two non-ideal configurations, distinguished by their \Gls{bsm} arm count, at $e_d=0.01$, $E_{\max}=100$, and $\gamma=0$.

\begin{table}[h]
\centering
\caption{Optimizer versus the paper schedule under two non-ideal network configurations. Both use $N=10$, $e_d=0.01$, $E_{\max}=100$, $\gamma=0$, and non-zero inner-qubit error $p_X^{\mathrm{in}}=p_Z^{\mathrm{in}}=3\times10^{-3}$, unlike the source paper's idealized inner-qubit assumption. The configurations differ in the number of redundant \Gls{bsm} arms per hop. A schedule is feasible when $F\ge F_{\min}=0.9$.}
\label{tab:netgen}
\small
\begin{tabular}{llcccc}
\hline
Config & Schedule & $C$ & $F$ & Rate & Feasible \\
\hline
18 arms & Paper schedule & 100 & 0.67 & 103.18 & \textbf{no} \\
18 arms & Matched-cost & 100 & 0.90 & 25.31 & yes \\
18 arms & Budget-relaxed & 80 & 0.90 & 85.30 & yes \\
\hline
6 arms & Paper schedule & 100 & 0.87 & 871.55 & \textbf{no} \\
6 arms & Matched-cost & 100 & 0.93 & 681.56 & yes \\
6 arms & Budget-relaxed & 60 & 0.91 & 2096.91 & yes \\
\hline
\end{tabular}
\end{table}

\begin{figure}[ht]
    \centering
    \includegraphics[width=0.85\textwidth]{figures/results/fidelity_by_config.png}
    \caption{Fidelity by network configuration and schedule. Hatched bars fall below the $F_{\min}=0.9$ threshold: the paper schedule is infeasible under both non-ideal configurations, while the schedules found by search remain feasible in every configuration tested.}
    \label{fig:netgen}
\end{figure}

The main result is a feasibility difference rather than a rate improvement: the paper schedule is infeasible in both non-ideal configurations, while the searched schedules remain feasible at the same cost. Raw rate comparisons are therefore not meaningful when the reference schedule fails the fidelity requirement.

The tree branching factor is not used as a sensitivity axis in this chapter because it does not affect $F$, $C$, or $R$ in the model of Chapter~\ref{chap:model}. Hop length is considered separately below; uniformly rescaling every hop length would multiply all schedule latencies by the same factor and leave relative schedule comparisons unchanged, whereas the heterogeneous hop lengths tested in Section~\ref{sec:exp:heterogeneous} change the relative physical conditions between hops.

\subsection{Per-Hop Heterogeneous Networks}
\label{sec:exp:heterogeneous}

% Source: outputs/random_network_adaptation/{30 seeds,100 seeds}/README.md,
%         src/hrgs_scheduler/models/random_network.py

Section~\ref{sec:exp:netgen} changes the noise model while keeping it uniform across hops. A real-world deployment, by contrast, can have different link lengths, hardware quality, and error rates from one hop to another. This section tests that setting by assigning each hop's length, inner-qubit error rates, and arm count independently at random. The search remains non-adaptive: the schedule is fixed once selected, and only the network configuration supplied to the search changes.

A five-hop network with one hop assigned a $12\times$ higher inner-qubit error rate than its neighbours provides a concrete test. At $E_{\max}=50$ and $F_{\min}=0.9$, the best uniform link-level candidate has $C=40$ and $F=0.88$ and is infeasible, while the optimizer finds a schedule with $C=38$, $F=0.90$, and $R=1798.57$ that clears the threshold by allocating resources differently across the hops.

\begin{table}[ht]
    \centering
    \small
    \begin{tabular}{lcccc}
        \hline
        Schedule & $C$ & $F$ & Rate & Feasible \\
        \hline
        Best uniform link-level recipe & 40 & 0.88 & 1738.32 & \textbf{no} \\
        Optimizer's schedule & 38 & 0.90 & 1798.57 & yes \\
        \hline
    \end{tabular}

    \caption{Designed weak-link contrast for a heterogeneous $N=5$ network at $e_d=0.01$, $E_{\max}=50$, and $F_{\min}=0.9$. One hop has an inner-qubit error rate $12\times$ higher than its neighbours. The uniform link-level baseline uses the same purification configuration at every hop, whereas the optimizer may allocate resources differently across hops. A schedule is feasible when $F\ge F_{\min}$.}
    \label{tab:weaklink}
\end{table}

The same tendency is observed over 100 independently generated random $N=6$ networks (seeds 0--99, $e_d=0.01$, $E_{\max}=60$, $F_{\min}=0.9$). Among the networks for which a feasible uniform link-level candidate exists, the optimizer improves rate by a mean of $20.37\%$ and never performs worse in the tested sample. Further interpretation of this behaviour is deferred to Chapter~\ref{chap:discussion}.

\subsection{Noise Sensitivity}
\label{sec:exp:edsweep}

% Source: outputs/sweep_ed_n10/README.md

Sweeping $e_d$ from 0 to 0.01 at $N=10$ shows that the optimizer's rate advantage over the uniform link-level baseline increases with noise, from near zero at $e_d=0$ to $18.2\%$ at $e_d=0.01$. Figure~\ref{fig:edsweep} shows the corresponding rate and fidelity curves for the paper baseline, the matched-cost optimizer, and the budget-relaxed optimizer. At low noise, the three variants remain close; the rate curves begin to separate visibly around $e_d=0.004$. The matched-cost optimizer may also change schedule structure as the noise level varies, so its curve should be interpreted as the best cost-$100$ result at each point rather than as one fixed schedule followed throughout the sweep. Further interpretation of these transitions is given in Chapter~\ref{chap:discussion}.

\begin{figure*}[ht]
    \centering
    \begin{subfigure}[b]{0.48\textwidth}
        \centering
        \includegraphics[width=\textwidth]{figures/results/rate_vs_ed.png}
        \caption{Rate vs.\ $e_d$}
    \end{subfigure}
    \hfill
    \begin{subfigure}[b]{0.48\textwidth}
        \centering
        \includegraphics[width=\textwidth]{figures/results/fidelity_vs_ed.png}
        \caption{Fidelity vs.\ $e_d$}
    \end{subfigure}
    \caption{Rate and fidelity vs.\ depolarizing error $e_d$ at $N=10$, $E_{\max}=100$, for the paper baseline, the optimizer at matched cost, and the budget-relaxed optimizer.}
    \label{fig:edsweep}
\end{figure*}

\section{Cost-Quality Tradeoffs}
\label{sec:exp:tradeoffs}

\subsection{Pareto Frontiers}
\label{sec:exp:pareto}

% Source: outputs/pareto_frontiers/README.md

With a permissive objective ($F_{\min}=0$, $E_{\max}=150$), the beam search evaluates 180 candidates. Figure~\ref{fig:pareto} shows two projections of the resulting Pareto frontier: fidelity versus cost and fidelity versus rate. Each point represents one evaluated schedule, while grey points are dominated candidates and connected points form the non-dominated frontier under the corresponding pair of objectives.

In the fidelity-cost projection, the paper schedule lies inside the frontier: a schedule at the same cost reaches approximately $0.006$ higher fidelity. The frontier is already close to flat beyond $C=60$, so additional cost provides only limited further fidelity improvement at this configuration. The paper schedule and the uniform link-level baseline are also marked for reference.

Appendix~\ref{app:dags} discusses the schedule \glspl{dag} corresponding to the key optima cited here and in Section~\ref{sec:exp:headline}.

\begin{figure*}[ht]
    \centering
    \begin{subfigure}[b]{0.48\textwidth}
        \centering
        \includegraphics[width=\textwidth]{figures/results/fidelity_vs_cost_n10.png}
        \caption{Fidelity vs.\ cost}
    \end{subfigure}
    \hfill
    \begin{subfigure}[b]{0.48\textwidth}
        \centering
        \includegraphics[width=\textwidth]{figures/results/fidelity_vs_rate_n10.png}
        \caption{Fidelity vs.\ rate}
    \end{subfigure}
    \caption{Pareto frontiers at $N=10$, $e_d=0.01$, and $E_{\max}=150$, over all 180 candidates evaluated by the beam search under a permissive objective ($F_{\min}=0$). Grey points are dominated candidates; connected points form the non-dominated frontier under each pair of objectives.}
    \label{fig:pareto}
\end{figure*}

\paragraph{Alternative objectives}
\label{sec:exp:objectives}

Changing only the optimization objective, from maximizing rate to maximizing fidelity or minimizing cost, produces different top-ranked schedules without changing the search algorithm. Maximizing fidelity subject to the paper's rate floor reaches $F=0.94$ at $C=80$, compared with $F=0.91$ at $C=60$ for the rate-maximizing objective. Minimizing cost subject only to $F\ge0.9$ returns $C=60$ but a rate as low as $R=1239.19$; adding an explicit rate floor restores $R=6195.95$ at the same cost. Appendix~\ref{app:objectives} gives the complete five-objective comparison.

\section{Scaling with Chain Length}
\label{sec:exp:scaling}

\subsection{Budget Scaling}
\label{sec:exp:budget}

% Source: outputs/sweep_min_budget_vs_n/README.md

For $N$ ranging from 10 to 28, the minimum $E_{\max}$ for which a schedule within the searched families reaches $F_{\min}=0.9$ is found by bisection, rerunning the beam search at each trial budget. A descriptive power-law fit to the resulting points in log-log space gives
\[
    E_{\max}^{\min} \approx 1.192\,N^{1.666},
\]
with an exponent greater than one over the measured range. The paper's budget is first exceeded at $N=28$, where the minimum feasible budget is $E_{\max}=283$, compared with $10\times28=280$; at $N=26$, the minimum is $258$, compared with $10\times26=260$. The observed crossover therefore lies between $N=26$ and $N=28$. Figure~\ref{fig:budget} shows the measured points together with the paper's linear formula and the fitted power law.

These fitted values describe the minimum budget found within the searched schedule families at the chosen beam width, not a global lower bound over the full legal schedule space or an asymptotic result (Section~\ref{sec:model:scope}). Chapter~\ref{chap:discussion} discusses this scope limitation together with the excluded-construction result of Section~\ref{sec:exp:excluded}.

Appendix~\ref{app:sweeps} additionally reports the minimum feasible budget as a function of $e_d$ at fixed $N=10$.

\begin{figure*}[ht]
    \centering
    \includegraphics[width=0.75\textwidth]{figures/results/min_budget_vs_n.png}
    \caption{Minimum feasible $E_{\max}$ required to reach $F_{\min}=0.9$ as a function of $N$ at $e_d=0.01$, compared with the paper's linear $E_{\max}=10N$ budget and the descriptive power-law fit $E_{\max}^{\min}\approx1.192\,N^{1.666}$. Solid markers and the solid portion of the fit cover the measured range ($N=10$--$28$); the dotted extension is the same fit evaluated beyond that range and does not represent additional measured data.}
    \label{fig:budget}
\end{figure*}

Appendix~\ref{app:results} (Section~\ref{sec:exp:excluded}) reports a related optimality-scope check at $N=18$, where a hand-constructed schedule using a same-span construction excluded from the search reaches the paper's own budget.

\section{Timing Sensitivity}
\label{sec:exp:timing}

% Source: outputs/sweep_gamma_and_tau_emit/README.md

Memory dephasing and generation latency affect schedules through their timing structure, including the placement of \texttt{HeraldNode}s and \texttt{IdleNode}s (Section~\ref{sec:model:dag}), rather than through aggregate fidelity or resource cost alone. Sweeping the memory dephasing rate $\gamma$ leaves the raw chain and paper schedule unchanged over the tested range, while heralded end-node pumping falls toward the maximally mixed state as $\gamma$ increases because its sacrificial copies remain in memory during the round-trip heralding delays. The schedules selected by the optimizer remain flat because the search favours optimistic schedules that avoid those intermediate waits. Sweeping the generation latency $\tau_{\mathrm{emit}}$ reduces the rate of every schedule; the detailed curves are given in Appendix~\ref{app:sweeps}. Chapter~\ref{chap:discussion} interprets these timing effects in the broader context of the scheduling problem.